% Part: sets-functions-relations
% Chapter: sets
% Section: subsets

\documentclass[../../../include/open-logic-section]{subfiles}

\begin{document}

\olfileid{sfr}{set}{sub}
\olsection{ઉપગણો અને ઘાતગણો}

\begin{explain}
આપણે ઘણી વાર ગણોની તુલના કરવા માગીશું. આવી તુલનાનો એક સ્પષ્ટ પ્રકાર આ છે:
\emph{એક ગણમાં જે કંઈ છે તે બધું બીજા ગણમાં પણ છે}.
આ પરિસ્થિતિ એટલી મહત્ત્વની છે કે તેના માટે આપણે નવી સંજ્ઞા રજૂ કરીએ છીએ.
\end{explain}

\begin{defn}[ઉપગણ]
જો ગણ $A$નો દરેક !!{element} $B$નો પણ !!a{element} હોય, તો આપણે કહીએ
કે $A$ એ $B$નો \emph{ઉપગણ} છે, અને $A \subseteq B$ લખીએ.
જો $A$ એ $B$નો ઉપગણ ન હોય, તો $A \not\subseteq B$ લખીએ.
જો $A \subseteq B$ પરંતુ $A \neq B$ હોય, તો $A \subsetneq B$ લખીએ
અને કહીએ કે $A$ એ $B$નો \emph{ઉચિત ઉપગણ} છે.
\end{defn}

\begin{ex}
દરેક ગણ પોતાનો જ ઉપગણ છે, અને $\emptyset$ દરેક ગણનો ઉપગણ છે.
યુગ્મ પ્રાકૃતિક સંખ્યાઓનો ગણ પ્રાકૃતિક સંખ્યાઓના ગણનો ઉપગણ છે.
વળી, $\{ a, b \} \subseteq \{ a, b, c \}$.
પરંતુ $\{ a, b, e \}$ એ $\{ a, b, c \}$નો ઉપગણ નથી.
\end{ex}

\begin{ex}
સંખ્યા $2$ પૂર્ણાંકોના ગણનો !!{element} છે, જ્યારે યુગ્મ સંખ્યાઓનો ગણ
પૂર્ણાંકોના ગણનો ઉપગણ છે. તેમ છતાં, કોઈ ગણ બીજા ગણનો !!a{element}
અને ઉપગણ \emph{બંને} હોઈ શકે. દા.ત.,
$\{0\} \in \{0, \{0\}\}$ અને $\{0\} \subseteq \{0, \{0\}\}$ પણ છે.
\end{ex}

ઘટકો દ્વારા નિર્ધારિત સમાનતાનો સિદ્ધાંત ગણો એક જ હોવાની કસોટી આપે છે:
$A = B$ ત્યારે અને ત્યારે જ થાય જ્યારે $A$નો દરેક !!{element}
$B$નો પણ !!a{element} હોય અને ઊલટું પણ સાચું હોય.
“ઉપગણ”ની વ્યાખ્યા અનુસાર $A \subseteq B$ એ આ કસોટીના પહેલા ભાગનું જ
નિરૂપણ છે: $A$નો દરેક !!{element} $B$નો પણ !!a{element} છે.
અલબત્ત, $A$ અને $B$ની અદલાબદલી કરીએ તો પણ આ વ્યાખ્યા લાગુ પડે છે:
$B \subseteq A$ ત્યારે અને ત્યારે જ થાય જ્યારે $B$નો દરેક !!{element}
$A$નો પણ !!a{element} હોય. આ તો ઘટકો દ્વારા નિર્ધારિત સમાનતાના
સિદ્ધાંતનો “ઊલટું પણ” એવો ભાગ જ છે.
બીજા શબ્દોમાં, આ સિદ્ધાંત અનુસાર ગણો ત્યારે અને ત્યારે જ સમાન હોય
જ્યારે તેઓ એકબીજાના ઉપગણો હોય.

\begin{prop}
$A = B$ ત્યારે અને ત્યારે જ થાય જ્યારે $A \subseteq B$ અને $B \subseteq A$ બંને હોય.
\end{prop}

કેટલીક વધુ ઉપયોગી સંજ્ઞાઓ રજૂ કરવાની આ સારી તક છે.
$A$ એ $B$નો ઉપગણ ક્યારે કહેવાય તેની વ્યાખ્યા આપતી વખતે આપણે
“$A$નો દરેક !!{element} \dots” કહ્યું હતું, અને “$\dots$”ની જગ્યાએ
“$B$નો !!a{element} છે” મૂક્યું હતું.
પરંતુ આવાં વિધાનોનું આ \emph{સ્વરૂપ} એટલું સામાન્ય છે કે તેના માટે
ઔપચારિક સંજ્ઞા રજૂ કરવી ઉપયોગી થશે.

\begin{defn}\ollabel{forallxina}
$(\forall x \in A)\phi$ એ $\forall x(x \in A \lif
\phi)$નું સંક્ષિપ્ત સ્વરૂપ છે. તેવી જ રીતે, $(\exists x \in A)\phi$
એ $\exists x(x \in A \land \phi)$નું સંક્ષિપ્ત સ્વરૂપ છે.
\end{defn}

આ સંજ્ઞા વાપરીને આપણે કહી શકીએ કે $A \subseteq B$ ત્યારે અને ત્યારે જ
થાય જ્યારે $(\forall x \in A)x \in B$ હોય.

હવે આપણે એક ખાસ પ્રકારના ગણનો વિચાર કરીએ: આપેલા ગણના બધા ઉપગણોનો ગણ.

\begin{defn}[ઘાતગણ]
ગણ~$A$ના બધા ઉપગણોથી બનેલા ગણને $A$નો \emph{ઘાતગણ} કહે છે
અને $\Pow{A}$ વડે દર્શાવીએ છીએ.
  \[
    \Pow{A} = \Setabs{B}{B \subseteq A}
  \]
\end{defn}

\begin{ex}
$\{ a, b, c \}$ના બધા શક્ય ઉપગણો કયા છે? તે આ છે:
$\emptyset$, $\{a \}$, $\{b\}$, $\{c\}$, $\{a, b\}$, $\{a, c\}$, $\{b,
c\}$, $\{a, b, c\}$. આ બધા ઉપગણોનો ગણ $\Pow{\{a,b,c\}}$ છે:
\[
\Pow{\{ a, b, c \}} = \{\emptyset, \{a \}, \{b\}, \{c\}, \{a, b\},
\{b, c\}, \{a, c\}, \{a, b, c\}\}
\]
\end{ex}

\begin{prob}
$\{a, b, c, d\}$ના બધા ઉપગણોની યાદી બનાવો.
\end{prob}

\begin{prob}
જો $A$માં $n$ !!{element}s હોય, તો $\Pow{A}$માં $2^n$
!!{element}s હોય તે બતાવો.
\end{prob}

\end{document}

